\section{ALEPH detector}
\label{sec:aleph}

ALEPH was a general-purpose LEP detector with near $4\pi$ coverage and high-precision charged-particle tracking in a 1.5~T solenoidal magnetic field~\cite{Decamp:1990jra,ALEPH:1994ayc}. The detector subsystems most relevant for this event-shape analysis are the central tracker, electromagnetic and hadronic calorimeters, and outer muon chambers.

The central tracking system combined a silicon vertex detector (VDET), an inner tracking chamber (ITC), and a large time projection chamber (TPC) with 21 pad rows covering $|\cos\theta|<0.97$. Together these detectors provided precise track parameters and momentum measurements, and they define the object-quality observables used in this analysis (for example $d_0$, $z_0$, and $N_{\mathrm{TPC}}$, defined in Sec.~\ref{sec:object_selection}). The single-track momentum resolution achieved by the combined tracking system is
$\sigma(1/p_T)\simeq 6\times10^{-4}$~GeV$^{-1}$ from the curvature measurement, to which a
multiple-scattering contribution must be added; it is negligible at high momentum and
dominates at low $p_T$. The corresponding constant terms for the calorimeter energy
resolutions below are likewise omitted here.
% NEEDS AUTHOR INPUT before submission: Referee 2 asks for the multiple-scattering term
% (they quote dpt/pt = 5e-3) and for the constant terms of the ECAL and HCAL energy
% resolutions. The full expressions are in Decamp:1990jra / ALEPH:1994ayc, both already
% cited above, but we have not read those papers and will not transcribe numbers we cannot
% check. Take the coefficients from the detector papers and state the complete expressions.
% Mirrored in referee/reviewer2.md.

The calorimeter system comprised two main subsystems. The electromagnetic calorimeter (ECAL) achieved an energy resolution of $\sigma_E/E \simeq 0.18/\sqrt{E\,[\mathrm{GeV}]}$, providing the primary response to photons and electromagnetic showers and anchoring the neutral-energy scale for event-shape reconstruction. The hadronic calorimeter (HCAL) with resolution $\sigma_E/E \simeq 0.85/\sqrt{E\,[\mathrm{GeV}]}$ is essential for neutral-hadron energy flow and global visible-energy reconstruction. The muon system outside HCAL supports lepton categorization in the archived ALEPH object definitions.

A key feature of the ALEPH reconstruction is the energy-flow algorithm, which combines tracking and calorimetric information to form a set of energy-flow objects that represent both charged and neutral particle candidates in the event. Tracking information dominates for charged particles, where the momentum resolution is superior to the calorimeter energy resolution. Calorimetric information provides the primary input for neutral particles (photons and neutral hadrons), which leave no track in the TPC. The combined energy-flow reconstruction achieves an energy resolution of $\sigma(E) = (0.59\pm0.03)\sqrt{E/\text{GeV}} + (0.6\pm0.3)$~GeV for hadronic events~\cite{Heister:2003aj}. This energy-flow approach, used in the ALEPH 2004 published event-shape analysis~\cite{Heister:2003aj}, is the particle-level input adopted for the present analysis. The 1998 ALEPH event-shape analysis~\cite{Barate:1996fi} used charged particles only to minimize model dependence on neutral-energy simulation; our nominal analysis follows the 2004 approach using both charged and neutral energy-flow objects, with a charged-only variation retained as a cross-check.

In this analysis, we use the archived ALEPH reconstructed object representation directly, including per-particle category labels (\texttt{pwflag}) and kinematic variables. Event-shape observables are built from selected charged and neutral energy-flow objects, so detector performance in each reconstructed particle category is a direct input to thrust migration and unfolding stability.
